We start by introducing the problem and its notations.

\begin{problem}[Nonlinear constrained optimization problem]
    \textbf{Given} the objective function $ f : \mathbf R^n \to \mathbf R $, and the constraint functions $ g : \mathbf R^n \to \mathbf R^m $ and $ h : \mathbf h \mathbf R^n \to R^p $,

    \textbf{Minimize} the function $ f $

    \textbf{Under the constraints}
    \[
        g(x) \succeq 0
    \]
    \[
        h(x) = 0
    \]
\end{problem}

In other words,

\begin{problem}[Nonlinear constrained optimization problem]
    \textbf{Given} the objective function $ f : \mathbf R^n \to \mathbf R $, and the constraint functions $ g : \mathbf R^n \to \mathbf R^m $ and $ h : \mathbf h \mathbf R^n \to R^p $,

    \textbf{Minimize} the function $ f $

    \textbf{Under the constraints}
    \[
        \forall i \in \{ 1, ..., m \}, g_i(x) \leqslant  0
    \]
    \[
        \forall j \in \{ j, ..., p \}, h_j(x) = 0
    \]
\end{problem}

\begin{definition}
    The domain of the problem is
    \[
        \mathcal D = \mathrm{dom} f \cap \bigcap_{i=1}^m \mathrm{dom} g_i \cap \bigcap_{j=1}^p \mathrm{dom} h_j
    \]

    The constraints set, or \textit{feasible set}, is
    \[
        X = \{ x \in \mathcal D, g(x) \preceq 0, h(x) = 0 \}
    \]

    The elements of $ X $ are called the the \textit{feasible points}, or \textit{primal feasible points}.
\end{definition}

\subsection{The Lagrangian function}

\begin{definition}[Lagrangian function]
    We define the \textit{Lagrangian dual function} associated to our constrained problem is
    \[
        \begin{aligned}
            L : \mathcal D \times \mathbf R^m \times \mathbf R^p & \longrightarrow \mathbf R \\
            (x,\lambda,\nu) & \longmapsto f(x) + \lambda^\top g(x) + \nu^\top h(x)
        \end{aligned}    
    \]

    The vectors $ \lambda $ and $ \nu $ are called the dual variables or Lagrange nultiplier vectors associated with the problem. We can refer to $ \lambda_i $ as the \textit{Lagrange nultiplier associated with the $i$-th inequality constraint $ g_i(x) \leqslant 0 $}, and similarly refer to $ \nu_j $ as the \textit{Lagrange nultiplier associated with the $j$-th equality constraint $ h_j(x) = 0 $}.
\end{definition}

\subsection{The Lagrange dual function}

\begin{definition}[Dual function]
    The \textit{dual function} associated to our constrained problem is
    \[
        \begin{aligned}
            d : \mathbf R^m \times \mathbf R^p & \longrightarrow \mathbf R \\
            (\lambda, \nu) & \longmapsto \inf_{x \in \mathcal D} L(x,\lambda,\nu)
        \end{aligned}
    \]
\end{definition}

\begin{problem}[Dual problem]
    \textbf{Maximize} the dual function associated to the primal problem
    \[ 
        d : \mathbf R^m \times \mathbf R^p \to \mathbf R
    \]

    \textbf{Under the constraint}
    \[
        \lambda \succeq 0
    \]
\end{problem}

In other words : the dual problem asks to compute
\[
    \sup_{\substack{ \lambda \succeq 0 \\ \nu \in \mathbf R^p }} d(\lambda, \nu)
\]

